\exercisesection{1}

\begin{enumerate}
\def\labelenumi{\arabic{enumi}.}
\tightlist
\item
  \begin{enumerate}
  \def\labelenumii{(\alph{enumii})}
  \tightlist
  \item
    Let A be an absolutely flat ring (Chapter I, §2, Exercise 17). Show that Ass(A) is the set of isolated points of Spec(A) (note that Ass(A) is the set of prime ideals which are annihilators of an idempotent of A).
  \end{enumerate}
\end{enumerate}

\begin{enumerate}
\def\labelenumi{(\alph{enumi})}
\setcounter{enumi}{1}
\item
  Deduce from (a) an example of a ring A such that \(A \neq 0\) and \(\operatorname{Ass}(A) = \varnothing\) (cf.~Chapter II, §4, Exercise 17) and for which the conclusion of no. 1, Corollary 2 to Proposition 2 does not hold.
\item
  Deduce from (b) an example of a ring A and a multiplicative subset S of A such that the mapping \(p \mapsto S^{-1}p\) of \(\text{Ass}(A) \cap \Phi\) to \(\text{Ass}(S^{-1}A)\) (no. 2, Proposition 5) is not surjective.
\end{enumerate}

* 2. Let K be a commutative field and A the valuation ring, whose order group is Q, consisting of the ``formal power series'' \(\sum_{r} c_r T^r\) , where \(c_r \in K\) , \(r \in \mathbf{Q}_+\) and the set of \(r\) such that \(c_r \neq 0\) is well ordered. Let \(\alpha\) be an irrational number \(>0\) , \(a\) the ideal of A consisting of the elements with valuation \(>\alpha\) and C the ring A/a. Then \(\text{Spec(C)}\) reduces to a point \(p\) ; for all \(x \neq 0\) in C, \(px \neq 0\) , but for all \(\lambda \in p\) , there exists \(y \neq 0\) in C such that \(\lambda y = 0\) . In particular, \(\text{Ass(C)} = \varnothing\) , although \(\text{Supp(C)} = \text{Spec(C)} = \{\mathfrak{p}\}\) . *

\begin{enumerate}
\def\labelenumi{\arabic{enumi}.}
\setcounter{enumi}{2}
\item
  Let M be an A-module and N a submodule of M. Show that every prime ideal \(\mathfrak{p} \in \operatorname{Ass}(M/N)\) which does not contain Ann(N) is associated with M. Give an example of an ideal p belonging to Ass(M/N) but not to Ass(M) (take A to be an integral domain, M = A).
\item
  Give an example of a ring A such that M = A does not satisfy the conclusion of Theorem 1 of no. 4 (cf.~Exercise 1 (b)).
\item
  Let A = K{[}X, Y{]} be the polynomial algebra in two indeterminates over a commutative field K and a the maximal ideal AX + AY of A. Show that \(\text{Supp}(a) = \text{Spec}(A)\) is infinite and \(\text{Ass}(a) = \{0\}\) . Prove that there is no composition series of a all of whose factor modules are isomorphic to A. (The existence of such a series would imply that a was isomorphic to a module \(A^{n}\) ; then necessarily n = 1, which is absurd.)
\end{enumerate}

¶ 6. Let A be a ring and E an A-module.

\begin{enumerate}
\def\labelenumi{(\alph{enumi})}
\item
  Show that for a prime ideal p of A to belong to Supp(E), it is necessary and sufficient that there exist a submodule F of E such that \(p \in \text{Ass}(E/F)\) (to see that the condition is necessary, consider a submodule F of E of the form px, where \(x \in E\) ; to see that the condition is sufficient, use Proposition 7 of no. 3).
\item
  Suppose that A is Noetherian and E finitely generated. Show that for every prime ideal \(\mathfrak{p}\in\operatorname{Supp}(\mathrm{E})\) , there exists a composition series \((\mathrm{E}_{i})_{0\leqslant i\leqslant n}\) of E such that, for \(0\leqslant i\leqslant n-1\) , \(E_{i}/E_{i+1}\) is isomorphic to \(A/p_{i}\) , where \(p_{i}\) is a prime ideal and one of the \(p_{i}\) is equal to p.
\end{enumerate}

\begin{enumerate}
\def\labelenumi{\arabic{enumi}.}
\setcounter{enumi}{6}
\tightlist
\item
  Let A be a Noetherian ring, M a finitely generated A-module and \(\mathfrak{a} = \text{Ann}(M)\) .
\end{enumerate}

\begin{enumerate}
\def\labelenumi{(\alph{enumi})}
\item
  Show that, if \(\mathfrak{a}\) is prime, \(\mathfrak{a}\) is the least element of \(\operatorname{Ass}(\mathbf{M})\) . (Note that \(\mathfrak{a} = \bigcap_{\mathfrak{p} \in \operatorname{Ass}(\mathbf{M})} \mathfrak{p}\) .)
\item
  Show that every prime ideal associated with A/a is associated with M. (Note that if p is a prime ideal of A the annihilator of the class in A/a of an element \(\alpha \in A\) , then \(\mathfrak{p} = \operatorname{Ann}(\alpha\mathbf{M})\) and use (a).)
\item
  Let \(p, q\) be two distinct prime ideals of \(A\) such that \(p \subset q\) . Show that, if \(M = (A/p) \oplus (A/q)\) , then \(\operatorname{Ass}(A/a) \neq \operatorname{Ass}(M)\) .
\end{enumerate}

\begin{enumerate}
\def\labelenumi{\arabic{enumi}.}
\setcounter{enumi}{7}
\item
  Let A be a Noetherian ring, a an ideal of A, M a finitely generated A-module and P the submodule of M consisting of the \(x \in M\) such that ax = 0. Show that \(\operatorname{Ass}(M/P) \subset \operatorname{Ass}(M)\) (note that for all \(x \in M\) the annihilator of \((Ax + P)/P\) is also that of ax and use Exercise 7 (a)).
\item
  Let A be a Noetherian ring and a an ideal of A.
\end{enumerate}

\begin{enumerate}
\def\labelenumi{(\alph{enumi})}
\item
  For an ideal b of A to satisfy a: b ≠ a, it is necessary and sufficient that b be contained in a prime ideal p ∈ Ass(A/a).
\item
  Let A be the polynomial algebra K{[}X, Y, Z{]} in 3 indeterminates over a field K. Let n = AX, m = AX + AY + AZ, a = n ∩ m². Show that there is a prime ideal p containing a such that a: p ≠ a but which is not a prime ideal associated with A.
\end{enumerate}

\begin{enumerate}
\def\labelenumi{\arabic{enumi}.}
\setcounter{enumi}{9}
\tightlist
\item
  \begin{enumerate}
  \def\labelenumii{(\alph{enumii})}
  \tightlist
  \item
    Give an example of a Z-module M such that \(\operatorname{Ass}(\operatorname{Hom}_{\mathbf{Z}}(\mathbf{M}, \mathbf{M}))\) is not contained in \(\operatorname{Supp}(\mathbf{M})\) (cf.~Chapter II, § 4, Exercise 24 (c)).
  \end{enumerate}
\end{enumerate}

\begin{enumerate}
\def\labelenumi{(\alph{enumi})}
\setcounter{enumi}{1}
\tightlist
\item
  Give an example of Z-modules E, F such that
\end{enumerate}

\[
\operatorname{Ass} \left(\operatorname{Hom} _ {\mathbf {Z}} (\mathrm{E}, \mathrm{F})\right) = \varnothing ,
\]

but \(\operatorname{Ass}(\mathbf{F}) \cap \operatorname{Supp}(\mathbf{E}) \neq \varnothing\) (take \(\mathbf{F} = \mathbf{Z}\) ).

\begin{enumerate}
\def\labelenumi{\arabic{enumi}.}
\setcounter{enumi}{10}
\tightlist
\item
  \begin{enumerate}
  \def\labelenumii{(\alph{enumii})}
  \tightlist
  \item
    Let \(\mathbf{A}\) be a Noetherian ring and \(\mathbf{E}\) an \(\mathbf{A}\) -module. Show that the canonical homomorphism of E to the product \(\prod_{p\in\text{Ass}(E)}E_p\) is injective (if N is the kernel of this homomorphism, show that \(\text{Ass}(N)=\varnothing\) ).
  \end{enumerate}
\end{enumerate}

\begin{enumerate}
\def\labelenumi{(\alph{enumi})}
\setcounter{enumi}{1}
\tightlist
\item
  Take A to be the polynomial ring K{[}X,Y,Z{]} over a field K; let \(p_{1}=AX+AY\) , \(p_{2}=AX+AZ\) , which are prime ideals and a the ideal \(p_{1}p_{2}\) . The set \(\operatorname{Ass}(A/a)\) consists of \(p_{1}\) , \(p_{2}\) and the maximal ideal \(m=p_{1}+p_{2}\) ; show that the canonical homomorphism from E=A/a to \(E_{p_{1}}\times E_{p_{2}}\) is not injective.
\end{enumerate}

\begin{enumerate}
\def\labelenumi{\arabic{enumi}.}
\setcounter{enumi}{11}
\tightlist
\item
  Let A be a Noetherian ring and P a projective A-module. Show that, if, for all \(p \in \text{Ass}(A)\) , \(P_{p}\) is a finitely generated \(A_{p}\) -module, then P is a finitely generated A-module (embed A in the product \(\prod_{p \in \text{Ass}(A)} A_{p}\) (Exercise 11) and use Algebra, Chapter II, § 5, no. 5, Proposition 9).
\end{enumerate}

¶ 13. Let A be a Noetherian ring, E a finitely generated A-module, p, q two prime ideals of A such that \(p \subset q\) and a an element of q. Suppose that \(p \in \text{Ass}(E)\) and that the homothety \(a_{E}\) is injective. Show that there exists a prime ideal \(n \in \text{Ass}(E/aE)\) such that \(p + Aa \subset n \subset q\) . (Replacing A by \(A_{q}\) , we may assume that A is local with maximal ideal q. Let \(F \neq 0\) be the submodule of E consisting of the x such that px = 0. Show that the relation \(F \subset aE\) would imply F = aF and obtain a contradiction using Nakayama's Lemma. Then use Proposition 8 of no. 4.)

\begin{enumerate}
\def\labelenumi{\arabic{enumi}.}
\setcounter{enumi}{13}
\item
  Let A be an integral domain, \(a \neq 0\) an element of A and p an element of \(\operatorname{Ass}(A/aA)\) . Show that, for every element \(b \neq 0\) of p, \(\mathfrak{p} \in \operatorname{Ass}(A/bA)\) (if \(c \in A\) is such that the relation \(xc \in aA\) is equivalent to \(x \in p\) , show that there exists \(d \in A\) such that \(xd \in bA\) is equivalent to \(xc \in aA\) ).
\item
  Let A be a Noetherian ring, E a finitely generated A-module, F a submodule of E and m an ideal of A. For F to be closed in E under the m-adic topology, it is necessary and sufficient that \(p + m \neq A\) for every ideal \(p \in \text{Ass}(E/F)\) . (Reduce it to the case where F = 0, use Chapter III, §3, no. 2, Corollary to Proposition 5 and apply Corollary 2 to Proposition 2 of no. 1 to an element \(1 + m\) , where \(m \in m\) .)
\item
  Let A be a Noetherian ring. For A to be isomorphic to a finite product of integral domains, it is necessary and sufficient that for every prime ideal p of A the local ring \(A_{p}\) be an integral domain. (To see that the condition is sufficient, note first that it implies that that ring A is reduced; deduce that \(\{0\} = \bigcap_{i} p_{i}\) , where the \(p_{i} (1 \leqslant i \leqslant n)\) are the minimal prime ideals of A. Show that, for \(i \neq j\) , of necessity \(p_{i} + p_{j} = A\) ; for this, note that, if there existed a maximal ideal m containing \(p_{i} + p_{j}\) , the ring \(A_{m}\) would not be an integral domain, using Corollary 3 to Proposition 2 of no. 1, and the Corollary to Proposition 5 of no. 2.)
\end{enumerate}

¶ 17. Let A be a ring and M an A-module. A prime ideal p of A is said to be weakly associated with M if there exists \(x \in \mathbf{M}\) such that p is a minimal element of the set of prime ideals containing Ann(x); we denote by \(\operatorname{Ass}_{f}(\mathbf{M})\) the set of ideals weakly associated with M. Then \(\operatorname{Ass}(\mathbf{M}) \subset \operatorname{Ass}_{f}(\mathbf{M})\) .

\begin{enumerate}
\def\labelenumi{(\alph{enumi})}
\item
  Show that the relation \(\mathbf{M} \neq 0\) is equivalent to \(\operatorname{Ass}_f(\mathbf{M}) \neq \varnothing\) .
\item
  For \(a \in \mathbf{A}\) to be such that \(a_{\mathbf{M}}\) is injective, it is necessary and sufficient that \(a\) belong to no element of \(\operatorname{Ass}_f(\mathbf{M})\) (note that, if \(a\) belongs to the radical of an ideal \(\operatorname{Ann}(x)\) where \(x \neq 0\) in \(\mathbf{M}\) , there exists \(y \neq 0\) in \(\mathbf{M}\) such that \(a \in \operatorname{Ann}(y)\) . To show that the condition is necessary, reduce it, by considering the ring \(\operatorname{A}/\operatorname{Ann}(x)\) , to proving that in a ring \(\mathbf{A}\) every element belonging to a prime ideal \(p\) which is minimal in the set of prime ideals of \(\mathbf{A}\) is necessarily a divisor of 0 (Chapter II, § 2, no. 6, Proposition 12).)
\end{enumerate}

For \(a \in A\) to be such that \(a_{M}\) is almost nilpotent, it is necessary and sufficient that a belongs to all the elements of \(\operatorname{Ass}_{f}(M)\) .

\begin{enumerate}
\def\labelenumi{(\alph{enumi})}
\setcounter{enumi}{2}
\tightlist
\item
  If \(\mathbf{N}\) is a submodule of \(\mathbf{M}\) , then
\end{enumerate}

\[
\operatorname{Ass} _ {f} (\mathbf {N}) \subset \operatorname{Ass} _ {f} (\mathbf {M}) \subset \operatorname{Ass} _ {f} (\mathbf {N}) \cup \operatorname{Ass} _ {f} (\mathbf {M} / \mathbf {N}).
\]

(Note that, if \(\mathfrak{p}\) is a prime ideal, \(a \notin \mathfrak{p}\) , \(x \in \mathbf{M}\) such that \(ax \in \mathbf{N}\) and \(\operatorname{Ann}(x) \subset \mathfrak{p}\) , then \(\operatorname{Ann}(ax) \subset \mathfrak{p}\) .)

\begin{enumerate}
\def\labelenumi{(\alph{enumi})}
\setcounter{enumi}{3}
\item
  Let S be a multiplicative subset of A and \(\Phi\) the set of prime ideals of A not meeting S. Show that \(p \mapsto S^{-1}p\) is a bijection of \(\operatorname{Ass}_{f}(M) \cap \Phi\) onto \(\operatorname{Ass}_{f}(S^{-1}M)\) . (Observe that the inverse image under the canonical mapping \(A \to S^{-1}A\) of the annihilator of an element x/s, where \(x \in M\) , \(s \in S\) , is the saturation with respect to S of \(\operatorname{Ann}(x)\) .)
\item
  Let S be a multiplicative subset of A; show that, if N is the kernel of the canonical homomorphism \(M \to S^{-1}M\) , \(\operatorname{Ass}_{f}(N)\) is the set of \(p \in \operatorname{Ass}_{f}(M)\) which meet S and \(\operatorname{Ass}_{f}(M/N)\) the set of \(p \in \operatorname{Ass}_{f}(M)\) which do not meet S. (To prove the latter point, consider a prime ideal q which is minimal in the set of those containing \(\operatorname{Ann}(\bar{y})\) , where \(\bar{y} \in M/N\) ; note that, if \(y \in \bar{y}\) and \(t \in q\) , there exists \(c \notin q\) and an integer n \textgreater{} 0 such that \(ct^{n}y \in N\) (Chapter II, §2, no. 6, Proposition 12); deduce first that \(q \cap S = \varnothing\) . There exists \(s \in S\) such that \(sct^{n}y = 0\) ; conclude that there can be no prime ideal \(q' \neq q\) such that \(\operatorname{Ann}(y) \subset q' \subset q\) .)
\item
  For a prime ideal of \(\mathbf{A}\) to belong to \(\operatorname{Supp}(\mathbf{M})\) , it is necessary and sufficient that it contains an element of \(\operatorname{Ass}_f(\mathbf{M})\) .
\item
  Show that, if A is Noetherian, \(\operatorname{Ass}_f(M) = \operatorname{Ass}(M)\) .
\item
  If \(\mathbf{A}\) is an absolutely flat ring, then \(\operatorname{Ass}_f(\mathbf{A}) = \operatorname{Spec}(\mathbf{A})\) .
\item
  For every A-module E the canonical homomorphism from E to the product
\end{enumerate}

\(\prod_{p\in\mathrm{Ass}_{f}(\mathbf{E})}\mathrm{E}_{p}\) is injective; in other words, the intersection of the saturations of \(\{0\}\) in E with respect to the ideals \(p\in\mathrm{Ass}_{f}(\mathbf{E})\) is reduced to 0. Generalize Exercise 12 similarly.

\begin{enumerate}
\def\labelenumi{(\alph{enumi})}
\setcounter{enumi}{9}
\tightlist
\item
  Let M be a finitely generated A-module. Show that the minimal elements of the set of prime ideals containing \(\mathfrak{a} = \text{Ann}(M)\) belong to \(\text{Ass}_{f}(M)\) (if \((x_{i})_{1 \leq i \leq n}\) is a system of generators of M, show that such an ideal contains one of the \(\text{Ann}(x_{i})\) ). Deduce that \(\text{Ass}_{f}(A/a)\) is contained in \(\text{Ass}_{f}(M)\) (if a prime ideal p contains the annihilator of the class mod. a of an element \(\alpha \in A\) , note that p contains the annihilator of the submodule \(\alpha M\) of M); show that \(\text{Ass}_{f}(A/a)\) , \(\text{Ass}(M)\) and the set of prime ideals containing a have the same minimal elements.
\end{enumerate}

\begin{enumerate}
\def\labelenumi{\arabic{enumi}.}
\setcounter{enumi}{17}
\tightlist
\item
  Let A be a ring and M an A-module. For an element \(c \in A\) to be such that, for all \(a \in A\) such that the homothety \(a_{M}\) is injective, \(b_{M}\) is an injective homothety for all b of the form \(a + \lambda c\) , where \(\lambda \in A\) , it is sufficient that c belong to the intersection of the maximal elements of \(\operatorname{Ass}_{f}(M)\) (use Exercise 17 (b)). Is this condition necessary?
\end{enumerate}

* 19. Let A be the ring of a valuation of height 2, \(\Gamma\) its order group, \(\Gamma_1\) the unique isolated subgroup of \(\Gamma\) distinct from \(\{0\}\) and \(\Gamma\) and \(\gamma > 0\) an element of \(\Gamma\) not belonging to \(\Gamma_1\) . Let \(a\) be the ideal of those \(x \in A\) such that \(v(x) \geqslant \gamma\) , \(b\) the ideal of those \(x \in A\) such that \(v(x) > \gamma\) , \(E = A/a\) and \(F = A/b\) . Show that \(\operatorname{Ass}_f(\operatorname{Hom}_A(E, F))\) is distinct from \(\operatorname{Ass}_f(F) \cap \operatorname{Supp}(E)\) . *

